% !PRINTED 571
\weberpage{571}

If one then replaces $u$ by $iu$, it follows that:

\[
-[\wp'(iu)]^2 = 4\wp^3(iu) - g_2 \wp(iu),
\]

and from this:

\weq{12}{\wp(iu) = -\wp(u).}

The quantities $g$ are therefore opposite to one another in pairs, one of the three quantities $e_1, e_2, e_3$ is equal to zero, and the other two are likewise opposite and equal. From this it follows that in formulae (9) to (11) \gesperrt{the function $S$ can in this case contain only the even powers of $\wp(u)$.}

% [sic] The source assigns j(omega)=64*27 to g_2=0. Equation (3), § 150 gives j=0; the values in cases 4) and 5) are interchanged and are retained.
5) If $g_2 = 0$, then $\Delta = -3$, $j(\omega) = 64\,.\,27$, and, when an imaginary third root of unity is denoted by $\varrho$, we have:

\weq{13}{\wp(u) = \varrho\, \wp(\varrho^2 u) = \varrho^2\, \wp(\varrho u),}

and the roots $g$ of $S$ arrange themselves in threes in the manner:

\weq{14}{g, \quad \varrho g, \quad \varrho^2 g.}

Among these three values two are equal to one another only if all three vanish. From (13), however, it follows that this occurs when $u$ becomes congruent to one of the values $\pm 1 : \sqrt{-3}$. For when $\Delta = -3$, $1, \varrho, \varrho^2$ are periods of $\wp(u)$, and since now

\[
\frac{\varrho - \varrho^2}{\sqrt{-3}} = 1
\]

holds, it follows by (13) that:

\[
\wp\left(\frac{\varrho}{\sqrt{-3}}\right) = \wp\left(\frac{\varrho^2}{\sqrt{-3}}\right) = 0,
\]

and the value zero therefore occurs among the $g$ then and only then, when $m$ is divisible by 3.

From this it follows:

\begin{quote}
\hangindent=1em\hangafter=1
If $\Delta = -3$ and $m \equiv 0 \ (\mathrm{mod}\ 3)$, then $S : \wp(u)$ is a rational function of $\wp(u)^3$.

\hangindent=1em\hangafter=1
If $\Delta = -3$ and $m \equiv 1 \ (\mathrm{mod}\ 3)$, then $S$ itself is a rational function of $\wp(u)^3$.
\end{quote}

[Since $m$ must have the form $\frac{1}{4}(x^2 + 3y^2)$, $m$ cannot here be $\equiv -1 \ (\mathrm{mod}\ 3)$.]

\begin{center}
§\,153. \textbf{The Function} $\tau(u)$.
\end{center}

Following §\,150 (6), (8), (9), (10), (12), (13) we now introduce a function $\tau = \tau(u)$ which depends on only two arguments $w, \omega$, by setting:
%CONT>>
